% =====================================================================
% Chapter: Why We Need Not Quantize Gravitation
%
% Uses only standard packages (amsmath, amssymb).
% Cite keys are prefixed "qg:" so they will not collide with the book's
% other bibliography entries. The thebibliography block at the end can be
% moved into the book's global .bib/.bbl if you prefer.
% =====================================================================

\chapter{Why We Need Not Quantize Gravitation}
\label{ch:no-quantum-gravitation}

Among the sentences of modern theoretical physics that are repeated as
if they were facts, one of the most common is: \emph{gravity must be
quantized}. It appears in textbook introductions, in popular books and in
the first paragraph of research proposals. It is much more rarely
accompanied by an argument. This chapter takes the sentence seriously
enough to ask three questions: what are the arguments, to which world
(the real one or the ideal one) does each of them belong, and does any
of them have the force of necessity?

The answer defended here is modest in form and uncomfortable in content.
We do not claim that gravitation \emph{is} classical. We claim that
\emph{no argument presently on the table shows that a fundamentally quantum
description of the gravitational field is needed}, and that the arguments
usually given are either demands of the ideal world, or concern regimes
to which the real world gives no causal access, or prove something weaker
than what they are quoted for. The majority of the community holds the
opposite view; this chapter is the case for the dissent, and we flag
honestly where the individual arguments are contested.

%---------------------------------------------------------------------
\section{Three kinds of ``need''}
%---------------------------------------------------------------------

A theory can be called \emph{necessary} in three quite different senses.

\begin{description}
\item[Empirical need.] There is an observation, or an experiment that can
  actually be performed, whose outcome cannot be described without the
  new theory.
\item[Logical need.] The existing theories, applied to a situation that can
  actually arise, give a contradiction or an ill-defined prediction.
\item[Unificatory need.] We would like to have one framework instead of
  two, and the new theory is the obvious candidate.
\end{description}

In the vocabulary of this book, the first is a need of the \emph{real
world}: it rests on experimental support. The second is a need of the
\emph{ideal world}, and it counts physically only if the offending
situation can be reached from the real world through a chain of
operational, causal steps; a contradiction among idealizations
that no experiment or observer can ever reach is a fact about the
idealizations, not about nature. The third is a purely ideal-world
need: a wish for logical tidiness. It is a perfectly respectable
research motive, but it is not a reason for nature to comply.

Our thesis can now be stated in three parts.

\begin{enumerate}
\item There is no empirical need for a quantum gravitational field: no
  observed phenomenon requires it.
\item The logical inconsistencies that are cited all live either in the
  ideal world only (extrapolations beyond any causal access), or turn out
  to establish less than is claimed (for example, that general relativity
  is an effective theory, which nobody disputes).
\item The unificatory need is real as a motive, and is not a necessity.
\end{enumerate}

``Need not'' is not ``must not''. Section~\ref{sec:change-mind} lists the
observations that would convert our position into its opposite.

%---------------------------------------------------------------------
\section{Three things called ``quantum gravity''}
%---------------------------------------------------------------------

Much of the confusion is verbal. The phrase \emph{quantum gravity}
covers at least three different things, and the claim ``gravity must be
quantized'' is only interesting for the third.

\begin{enumerate}
\item[(a)] \textbf{Quantum fields on a classical curved spacetime.} Matter and
  radiation are quantum, the metric is a fixed classical background (or
  is classical and dynamical, but not itself quantized). This is a
  mature, well-tested framework. Particle creation in an expanding
  universe and Hawking radiation are computed here; so is the spectrum of
  primordial density perturbations in inflationary cosmology. None of
  these is a quantization of the metric.

\item[(b)] \textbf{General relativity as an effective field theory.} One
  \emph{can} quantize the metric perturbatively around flat space. The
  theory is non-renormalizable, but as an effective theory it is
  perfectly predictive at energies far below the Planck scale. Donoghue
  \cite{qg:donoghue} computed the leading quantum correction to the
  Newtonian potential between two masses $m_1$, $m_2$:
  \begin{equation}
    V(r) = -\frac{G m_1 m_2}{r}
    \left[\, 1 + 3\,\frac{G(m_1+m_2)}{r c^2}
      + \frac{41}{10\pi}\,\frac{G\hbar}{r^2 c^3} \right].
    \label{eq:donoghue}
  \end{equation}
  The second term is classical (a post-Newtonian correction); the third is
  the genuinely quantum one. It is of relative order
  $(\ell_P/r)^2$, where $\ell_P=\sqrt{\hbar G/c^3}\approx 1.6\times
  10^{-35}\,\mathrm{m}$. At $r\sim 10^{-19}\,\mathrm{m}$, roughly the
  shortest distance probed at present-day colliders, this is of order
  $10^{-32}$. So ``quantum gravity'' in sense (b) exists, is harmless, and
  is unobservable.

\item[(c)] \textbf{A fundamental, UV-complete quantum theory of the
  gravitational field (of spacetime geometry).} This is the programme
  usually meant by the slogan, and the one whose \emph{necessity} is at
  issue.
\end{enumerate}

Sense (a) is already in the textbook, sense (b) is a technical fact, and
neither is under dispute. Everything below concerns (c).

%---------------------------------------------------------------------
\section{The arguments for, one by one}
%---------------------------------------------------------------------

%.....................................................................
\subsection{``Everything else is quantum''}
%.....................................................................

The most common argument is an induction: electromagnetic, weak and
strong interactions have been brought into the quantum framework, so
gravitation must follow. Two remarks.

First, an induction over three cases is only as good as the
\emph{similarity} of the fourth. The three quantized interactions are
gauge theories of one type: spin-1 mediators and internal symmetries on a
fixed spacetime. Gravitation is of a different kind: it \emph{is} the
spacetime. Einstein's equations can be read as an equation of state
rather than as a fundamental field equation (Section~\ref{sec:emergent});
a thermodynamic or hydrodynamic equation is not usually considered a
candidate for quantization.

Second, what has actually been quantized in the three successful cases
is a field that carries \emph{energy and momentum through spacetime}. For
the gravitational field, the analogous step is anything but unambiguous:
the notion of local energy of the gravitational field is itself
frame-dependent, and ``the same'' physics can be expressed in terms of
metric, tetrads, connections or loops, which give inequivalent quantum
theories. The induction, in other words, has no unique conclusion to
arrive at.

An observational remark belongs here too. Gravity is feeble: the ratio of
the gravitational to the electrostatic force between two protons is
$G m_p^2 \big/ (e^2/4\pi\varepsilon_0)\approx 8\times 10^{-37}$. A
single graviton has never been detected, and Dyson \cite{qg:dyson} and
Rothman and Boughn \cite{qg:rothman} have argued that it very probably
cannot be, even in principle: a detector sensitive enough to register one
graviton would have to be so massive that it collapses into a black hole
first. These arguments are not proofs of classicality, but they show that
the quantum nature of the gravitational field, if it exists, will not be
verified by the usual route of counting quanta.

%.....................................................................
\subsection{Consistency of the coupling}
%.....................................................................

The most serious argument is a genuine one: how can a classical field be
coupled to quantum matter? The naive answer is the semiclassical
Einstein equation
\begin{equation}
  G_{\mu\nu} = \frac{8\pi G}{c^4}\,\langle \hat T_{\mu\nu}\rangle ,
  \label{eq:semiclassical}
\end{equation}
where the source is the quantum expectation value. In the non-relativistic
limit this becomes the Schr\"odinger--Newton equation
\cite{qg:diosi84,qg:penrose}
\begin{equation}
  i\hbar\,\partial_t \psi(\mathbf r,t) =
  \left[-\frac{\hbar^2}{2m}\nabla^2
    - G m^2 \!\int\! \frac{|\psi(\mathbf r',t)|^2}{|\mathbf r-\mathbf r'|}
    \,d^3r'\right]\psi(\mathbf r,t),
  \label{eq:sn}
\end{equation}
which is non-linear in $\psi$. The superposition principle fails, and
Eq.~\eqref{eq:semiclassical} is therefore regarded as unacceptable by many
authors \cite{qg:kibble}.

Two famous arguments are directed at this. Eppley and Hannah
\cite{qg:eppley} proposed a thought experiment in which a classical
gravitational wave is used to probe a quantum particle, and claimed that
either the uncertainty principle or momentum conservation must be
violated. Mattingly \cite{qg:mattingly} argued that the thought
experiment does not establish the claim; Albers, Kiefer and Reginatto
\cite{qg:albers}, using an analysis of measurability in the spirit of Bohr
and Rosenfeld, argued broadly the other way, namely that a classical
gravitational field coupled to quantum matter leads to inconsistencies
unless quantum mechanics itself is modified. Page and Geilker
\cite{qg:page} analysed a Cavendish-type arrangement in which the position
of the source mass depends on the outcome of a quantum measurement; the
field a real Cavendish apparatus would register follows the \emph{outcome},
not the expectation value. They concluded that either the wave function
really collapses, or the gravitational field is quantum.

This is the strongest card of the opposite side, and we should say what
it does and does not show.

\begin{itemize}
\item It shows that the \emph{naive} coupling \eqref{eq:semiclassical}
  cannot be the last word, especially for interpretations of quantum
  mechanics without collapse. This is a result about one particular
  coupling prescription, not about classical gravitation as such.
\item It does not show that no consistent classical--quantum coupling
  exists. Oppenheim \cite{qg:oppenheim} has constructed a fully
  consistent dynamics for quantum matter coupled to a \emph{classical}
  spacetime metric. The price is that the dynamics is stochastic: the
  classical metric fluctuates, and there is a trade-off between the
  decoherence of the quantum system and the diffusion of the classical
  field. This is not a defect: it is a testable prediction, and it
  identifies where to look.
\item Proposals in which gravity is the \emph{agent} of wave-function
  collapse \cite{qg:diosi84,qg:penrose} go in the same direction: they
  treat gravity as the reason why macroscopic superpositions do not
  persist, i.e.\ as the part of the world that does not obey the
  superposition principle.
\end{itemize}

Whether any of these succeeds is for experiment to decide. The point for
the present chapter is that the logical gap exists: ``classical gravity
coupled to quantum matter is inconsistent'' is not a theorem.

%.....................................................................
\subsection{Singularities}
%.....................................................................

General relativity predicts singularities under rather general
conditions \cite{qg:hawkingpenrose}, and quantum gravity is expected to
remove them. Three comments.

\emph{A singularity is a statement about the limit of applicability, and a
theory that announces the limit of its own applicability is behaving
well.} Geometrical optics predicts infinite intensity at a caustic. The
cure is wave optics, which is a different theory; it is not obtained by
``quantizing rays''. In the same way, a curvature singularity indicates
that something other than the Einstein equations governs the situation.
It does not say \emph{what}: it may be a modified classical gravitational
dynamics at high curvature, different matter equations of state, or
something structurally unlike a quantized metric. The theorems also
require energy conditions and global assumptions, which are
idealizations of matter.

Second, the singularities hide. Inside a black hole they are causally
disconnected from every outside observer, and the cosmic-censorship
conjecture expresses the hope that this is generic. Physics that cannot
causally influence any observer lies outside the real world in the sense of
this book.

Third, the Big Bang singularity is an extrapolation. Our observational
reach into the early universe ends at nucleosynthesis and, with some
indirect evidence, at an inflationary epoch. The extrapolation of a smooth
solution to $t=0$ is an exercise in the ideal world.

%.....................................................................
\subsection{Black holes and information}
%.....................................................................

Hawking's calculation \cite{qg:hawking} indicates that black holes
evaporate by thermal radiation and raises the question of whether
information is lost. The paradox is often presented as a proof that the
semiclassical description fails and quantum gravity is required. Consider
its ingredients: (i) a quantum-field-theory calculation on a classical
background (sense (a) above); (ii) the demand that the global evolution
be unitary; (iii) the event horizon, which is a \emph{teleological}
concept, as it can only be defined with the entire future of the
spacetime in hand. Local, operationally accessible notions (apparent
horizons) differ from it.

This is a typical \emph{mixture of the two worlds}. Item (ii) is a
requirement of the ideal-world formalism. Item (iii) is a construction
that depends on information that no real observer has, and cannot
have, at the time at which the question is posed. A paradox that
results from combining the three is a signal that one of the ideal-world
assumptions has been pushed too far, and which one is a question about
\emph{logic}, not about whether the metric is an operator.

The real-world side is less dramatic than the literature suggests. A
solar-mass black hole has a Hawking temperature of about
$6\times10^{-8}\,\mathrm{K}$, far below that of the cosmic microwave
background (2.7~K), so it currently absorbs more than it emits; its
evaporation time is of order $10^{67}$ years. No evaporating black hole
has been observed. Indeed, the Hawking effect requires nothing but
quantum matter in a curved classical background, and analogue systems have
reportedly displayed it \cite{qg:steinhauer}: such a signal does not test
quantum gravity at all.

%.....................................................................
\subsection{The Planck scale}
%.....................................................................

The Planck length $\ell_P\approx 1.6\times10^{-35}\,\mathrm{m}$ and Planck
energy $E_P=\sqrt{\hbar c^5/G}\approx 1.2\times10^{19}\,\mathrm{GeV}$ are
said to mark the point at which the quantum fluctuations of the metric
become of order one. This is the least empirical of the arguments. The
Planck energy is some fifteen orders of magnitude above the energies
reached at the LHC, and more than seven orders above the most
energetic cosmic rays, even in the laboratory frame.

Moreover, $\ell_P$ is built by dimensional analysis from $G$, $\hbar$,
$c$. Dimensional analysis can produce a length from any three dimensional
constants. Stoney did this before Planck with $G$, $c$ and $e$:
\begin{equation}
  \ell_S = \sqrt{\frac{G e^2}{4\pi\varepsilon_0 c^4}}
  \approx 1.4\times10^{-36}\,\mathrm{m}
  = \sqrt{\alpha}\;\ell_P ,
\end{equation}
and nobody concludes that something special happens at the Stoney length.
Dimensional analysis can tell us where an effective theory \emph{might}
cease to be valid; it does not tell us what takes over.

%.....................................................................
\subsection{Non-renormalizability}
%.....................................................................

The coupling constant of the perturbative quantization of general
relativity has negative mass dimension, $G\sim M_P^{-2}$, and the theory
is non-renormalizable. Fermi's four-fermion theory of weak interaction had
the same property, with $G_F\approx 1.2\times10^{-5}\,\mathrm{GeV}^{-2}$,
signalling a breakdown near $G_F^{-1/2}\sim 300\,\mathrm{GeV}$. The
resolution was the electroweak theory, with new particles $W^\pm$ and
$Z$. Nobody calls it ``quantum Fermi theory'': it is a different theory
from which Fermi's emerges at low energy.

We accept the lesson: non-renormalizability \emph{does} point to new
physics at some scale. What it does not point to is that the new physics is
a quantization of the metric. Indeed, if general relativity is
non-perturbatively renormalizable (``asymptotic safety''
\cite{qg:weinberg79}), it would mean that a UV-complete quantum theory of
the metric exists; it would still not follow that any real-world situation
requires it.

%---------------------------------------------------------------------
\section{Diagnosis: two maps, one territory}
\label{sec:maps}
%---------------------------------------------------------------------

Why does the demand for quantization feel so compelling? Because
general relativity and quantum mechanics are two \emph{ideal-world}
constructions that are formally incompatible: one has a dynamical
causal structure and no background; the other has linear evolution on a
fixed background and a distinguished time. The incompatibility is real
and has produced the ``problem of time'' and other conceptual puzzles.
But the incompatibility is between two mathematical models, and the
temptation is to read it as a contradiction in nature.

Recall the cartographic picture that runs through this book. A physical
theory is a map, not the territory; different theories are maps of
different kinds, with different projections and different resolutions.
Gauss's \emph{Theorema Egregium} states that a sphere cannot be mapped to
a plane without distortion, because the Gaussian curvatures differ.
Every projection preserves something (angles in Mercator's, areas in
Lambert's) and sacrifices something else. We do not conclude that the Earth
is inconsistent. We do not even demand a single flat map that preserves
everything. Quantum mechanics preserves linear superposition and
unitary evolution on a given background; general relativity preserves a
dynamical causal structure and background independence. To demand one map
that preserves all of these is a demand, and may be unreasonable as a
demand.

And a map does not have to be single in order to be mathematically
complete. A manifold is \emph{defined} by an atlas: charts with
transition rules between them. If the territory is best described by an
atlas of theories, the semiclassical coupling prescriptions of the previous
section play the role of transition functions, and the open task is to
state them correctly (the programme sketched by the classical--quantum
hybrid theories above), not to produce a single chart that covers
everything.

%---------------------------------------------------------------------
\section{Where a quantum theory might live, and why it need not be a
quantization of gravity}
\label{sec:emergent}
%---------------------------------------------------------------------

Suppose, against the thesis, that nature at its deepest level is quantum.
Does it follow that the Einstein equations must be quantized? No, and
the history of physics gives clear examples to the contrary.

Take a crystal. Its elastic continuum theory is an excellent
description at long wavelengths. One \emph{can} quantize the elastic field
and obtain phonons, and the result is correct at low energies. But
phonons are not fundamental, and the fundamental quantum theory is that of
the atoms. The quantization of the continuum theory does not lead to the
atoms; an observer who quantized elasticity and stopped there would
never discover them.

Gravitation may be like this. The idea that gravity is an emergent,
collective, thermodynamic phenomenon has a respectable pedigree:
Sakharov's induced gravity \cite{qg:sakharov}; Jacobson's derivation of the
Einstein equation from the thermodynamics of local Rindler horizons
\cite{qg:jacobson}, in which the equation is an \emph{equation of state};
and the extensive literature on analogue gravity, in which the
propagation of excitations in a moving medium obeys the equations of
a curved-space field theory \cite{qg:analogue}. In these examples a
metric appears, with causal cones and horizons, without anyone having
quantized a metric; the underlying quantum theory is that of the medium.

There is an honest caveat. The Weinberg--Witten theorem \cite{qg:ww}
forbids a massless spin-2 particle coupled to a Lorentz-covariant
conserved stress tensor in a Lorentz-covariant flat-space theory, so
emergent gravity cannot be \emph{too naive}: it needs a loophole such as
non-trivial background, emergent Lorentz symmetry, or holography.
Holography (AdS/CFT) is, however, a case in point rather than a
counter-example: the fundamental theory is a non-gravitational quantum
field theory on the boundary, and the bulk gravity is a derived
description. The quantum theory is real, and ``quantum gravity'' is its
name for some of its regimes.

The conclusion is that \emph{even if} a quantum theory underlies spacetime,
the correct object of study is that underlying theory, and the
quantization of the Einstein equations is the wrong target. The slogan
``gravity must be quantized'' confuses the two.

%---------------------------------------------------------------------
\section{What would change our mind}
\label{sec:change-mind}
%---------------------------------------------------------------------

A position of this kind must say in advance what would refute it. The
real-world criteria are as follows.

\begin{enumerate}
\item \textbf{Gravitationally mediated entanglement.} Bose et
  al.\ \cite{qg:bose} and Marletto and Vedral \cite{qg:marletto} proposed
  an experiment in which two mesoscopic masses (of order $10^{-14}$\,kg),
  each held in a spatial superposition of order $100\,\mu\mathrm{m}$ and
  coupled only through gravity, become entangled. If two systems become
  entangled through a mediator under local operations only, the mediator
  cannot be classical. Observation of such entanglement would create a
  real-world need for a non-classical gravitational degree of freedom,
  and our thesis would fall. The experiment has not been done at the time
  of writing.
\item \textbf{Fundamental gravitational noise.} The classical-field
  hypothesis of \cite{qg:oppenheim} predicts irreducible stochastic
  fluctuations and a lower bound on decoherence of massive systems tied
  to them. Its observation, and the absence of the corresponding
  quantum-gravitational entanglement, would support the classical
  hypothesis; failure to find either at the predicted level would
  constrain it.
\item \textbf{Primordial gravitational waves.} Krauss and Wilczek
  \cite{qg:krauss} argued that the detection of primordial tensor
  modes would be evidence for the quantization of gravity. The claim is
  disputed, because it is not clear that the same spectrum could not be
  produced by other mechanisms, but a detection would at least shift the
  burden of proof.
\item \textbf{Any breakdown of the effective theory that is accessible.}
  A deviation from Eq.~\eqref{eq:donoghue}-type predictions, or from
  Lorentz invariance, at accessible energies would be an immediate real-world
  need.
\end{enumerate}

None of these has produced a positive result so far, which is the
statement that there is currently no empirical need.

%---------------------------------------------------------------------
\section{Objections and replies}
%---------------------------------------------------------------------

\begin{itemize}
\item \textbf{``Hybrid theories are ad hoc.''} They are, at the moment, less
  developed than the quantum programme, and they pay for consistency with
  stochasticity. But ``less developed'' is not ``impossible'', and the
  construction of \cite{qg:oppenheim} shows that the impossibility claim
  is not a theorem.
\item \textbf{``The universe contains regimes where gravity and quantum effects are
  both strong.''} Yes: the interior of black holes, the first instants of the
  universe, and Planck-energy collisions. In each of them the real world
  provides no causal access, which is exactly why the arguments from these
  regimes are arguments of the ideal world.
\item \textbf{``Effective field theory shows that gravity can be quantized.''} It does,
  and we have said so (sense (b)). That is the reason why we need not
  quantize anything \emph{further} in order to describe what is observed.
\item \textbf{``You are assuming that nature is not unified.''} No. We say that
  unification is a legitimate aim and that its success does not
  \emph{need} to take the form ``quantize the metric''. If a unified theory
  exists, it may very well be a quantum theory from which gravitation
  emerges; this is compatible with everything above.
\end{itemize}

%---------------------------------------------------------------------
\section{Summary}
%---------------------------------------------------------------------

\begin{itemize}
\item The slogan ``gravity must be quantized'' mixes three meanings:
  quantum fields on a classical background (established), quantization of
  general relativity as an effective theory (consistent and unobservable),
  and a fundamental quantum theory of the metric (not shown to be needed).
\item The strongest argument, the inconsistency of classical gravity
  coupled to quantum matter, refutes a particular naive coupling, not
  classical gravitation; consistent classical--quantum hybrid dynamics
  exist and make testable predictions.
\item Singularities, the information paradox and the Planck scale are
  statements about extrapolations beyond causal access, i.e.\ about the
  ideal world; they signal that general relativity is an effective theory,
  not that its quantization is the successor.
\item The right analogy is a map, not a derivation: theories are
  charts of different projections, an atlas is a legitimate mathematical
  object, and the fact that two charts cannot be glued into one flat map
  is not an inconsistency of the territory.
\item If a quantum theory underlies spacetime, it is the theory of the
  underlying degrees of freedom and not the quantization of the Einstein
  equations: phonons are not atoms.
\item The position is falsifiable: observation of gravitationally mediated
  entanglement would create an empirical need for a quantum gravitational
  degree of freedom.
\end{itemize}

%---------------------------------------------------------------------
\section*{Notes and references}
%---------------------------------------------------------------------

\begin{thebibliography}{99}

\bibitem{qg:eppley} K.~Eppley and E.~Hannah, The necessity of quantizing
  the gravitational field, \emph{Found.\ Phys.} \textbf{7} (1977) 51--68.

\bibitem{qg:page} D.~N.~Page and C.~D.~Geilker, Indirect evidence for
  quantum gravity, \emph{Phys.\ Rev.\ Lett.} \textbf{47} (1981) 979--982.

\bibitem{qg:kibble} T.~W.~B.~Kibble, Is a semi-classical theory of
  gravity viable?, in \emph{Quantum Gravity 2: A Second Oxford Symposium},
  eds.\ C.~J.~Isham, R.~Penrose, D.~W.~Sciama (Oxford University Press,
  1981).

\bibitem{qg:mattingly} J.~Mattingly, Why Eppley and Hannah's thought
  experiment fails, \emph{Phys.\ Rev.\ D} \textbf{73} (2006) 064025.

\bibitem{qg:albers} M.~Albers, C.~Kiefer and M.~Reginatto, Measurement
  analysis and quantum gravity, \emph{Phys.\ Rev.\ D} \textbf{78} (2008)
  064051.

\bibitem{qg:carlip} S.~Carlip, Is quantum gravity necessary?,
  \emph{Class.\ Quantum Grav.} \textbf{25} (2008) 154010. (A general
  review of the question; close in spirit to this chapter.)

\bibitem{qg:dyson} F.~Dyson, Is a graviton detectable?,
  \emph{Int.\ J.\ Mod.\ Phys.\ A} \textbf{28} (2013) 1330041.

\bibitem{qg:rothman} T.~Rothman and S.~Boughn, Can gravitons be detected?,
  \emph{Found.\ Phys.} \textbf{36} (2006) 1801--1825.

\bibitem{qg:oppenheim} J.~Oppenheim, A postquantum theory of classical
  gravity?, \emph{Phys.\ Rev.\ X} \textbf{13} (2023) 041040.

\bibitem{qg:diosi84} L.~Di\'osi, Gravitation and quantum-mechanical
  localization of macro-objects, \emph{Phys.\ Lett.\ A} \textbf{105}
  (1984) 199--202.

\bibitem{qg:penrose} R.~Penrose, On gravity's role in quantum state
  reduction, \emph{Gen.\ Relativ.\ Gravit.} \textbf{28} (1996) 581--600.

\bibitem{qg:donoghue} J.~F.~Donoghue, General relativity as an effective
  field theory: the leading quantum corrections, \emph{Phys.\ Rev.\ D}
  \textbf{50} (1994) 3874--3888.

\bibitem{qg:hawkingpenrose} S.~W.~Hawking and R.~Penrose, The
  singularities of gravitational collapse and cosmology, \emph{Proc.\ R.\
  Soc.\ Lond.\ A} \textbf{314} (1970) 529--548.

\bibitem{qg:hawking} S.~W.~Hawking, Breakdown of predictability in
  gravitational collapse, \emph{Phys.\ Rev.\ D} \textbf{14} (1976)
  2460--2473.

\bibitem{qg:steinhauer} J.~Steinhauer, Observation of quantum Hawking
  radiation and its entanglement in an analogue black hole,
  \emph{Nature Phys.} \textbf{12} (2016) 959--965.

\bibitem{qg:weinberg79} S.~Weinberg, Ultraviolet divergences in quantum
  theories of gravitation, in \emph{General Relativity: An Einstein
  Centenary Survey}, eds.\ S.~W.~Hawking and W.~Israel (Cambridge
  University Press, 1979), pp.~790--831.

\bibitem{qg:sakharov} A.~D.~Sakharov, Vacuum quantum fluctuations in
  curved space and the theory of gravitation, \emph{Dokl.\ Akad.\ Nauk
  SSSR} \textbf{177} (1967) 70--71.

\bibitem{qg:jacobson} T.~Jacobson, Thermodynamics of spacetime: the
  Einstein equation of state, \emph{Phys.\ Rev.\ Lett.} \textbf{75}
  (1995) 1260--1263.

\bibitem{qg:analogue} C.~Barcel\'o, S.~Liberati and M.~Visser, Analogue
  gravity, \emph{Living Rev.\ Relativ.} \textbf{14} (2011) 3.

\bibitem{qg:ww} S.~Weinberg and E.~Witten, Limits on massless particles,
  \emph{Phys.\ Lett.\ B} \textbf{96} (1980) 59--62.

\bibitem{qg:bose} S.~Bose \emph{et al.}, Spin entanglement witness for
  quantum gravity, \emph{Phys.\ Rev.\ Lett.} \textbf{119} (2017) 240401.

\bibitem{qg:marletto} C.~Marletto and V.~Vedral, Gravitationally induced
  entanglement between two massive particles is sufficient evidence of
  quantum effects in gravity, \emph{Phys.\ Rev.\ Lett.} \textbf{119}
  (2017) 240402.

\bibitem{qg:krauss} L.~M.~Krauss and F.~Wilczek, Using cosmology to
  establish the quantization of gravity, \emph{Phys.\ Rev.\ D}
  \textbf{89} (2014) 047501.

\end{thebibliography}
